% Folder 136, batch 01 (archivists' pages 1-20).
% Pass: Opus 5.5 (claude-opus-5-5), 2026-10-03 — first pass, unchecked against the pages by a human.
%
% Page 1 is an orange cover inscribed in his hand; its words become the first
% section heading. Page 18 is a blank coloured sheet. Both carry no \page{}.
\documentclass[11pt,a4paper]{article}
\input{../preamble/grothendieck}

\folder{136}
\batch{1}
\pages{1}{20}
\foldertitle{Complexe de De Rham à puissance divisée [conférence de 1976 à l’IHÉS] : notes manuscrites (s.d.).}
\dating{[à partir de 1975-1976]}
\watermark{Édition de démonstration}

\begin{document}

\section{Complexe DR à p.d.\ et ombres}
\note{titre inscrit de sa main sur la couverture orange (p.~1) ; la série qui suit est numérotée par l'auteur de (1) à (16), pp.~2--17.}
\page{2}
\note{En tête de la page, numérotée (1) par l'auteur.}
\marginal{au crayon, encadré dans le coin supérieur gauche : Plan conf. IHÉS en 1976}
\marginal{au crayon, en haut à droite : géom.\ diff. / — anal. / — algébr. / \struck{\ill{}} arithm. / top.\ alg.\ $\{$PL, ½simpl.}

\textbf{Complexe DR à p.d.\ et ombres des modules.}

1 \quad Historique

a) Notion de forme diff.\ (\emph{Poincaré}) et formule de \emph{Stokes}
\[ \int_C d\omega = \int_{\partial C} \omega \]

b) Th.\ de \emph{De Rham} (conjecturé par \emph{E.~Cartan}) (dit $H^*_{\mathrm{DR}}(X) \to H^*(X,\mathbb{C})$). \uncertain{Mais on perd la torsion.} Maintenant bien compris : th.\ des faisceaux, lemme Poincaré.

c) Théorie de \emph{Hodge des} intégr.\ harmoniques (structure supplémentaire bigraduée sur $H^*_{\mathrm{DR}}(X)$ si $X$ kählérienne compacte)

d) Théorème de \emph{H.~Cartan}–\emph{Serre} sur variétés de Stein (où lemme de Poincaré valable dans l'analytique)

e) Cas des variétés \emph{algébriques} ou schémas, sur \add{corps de} base \ill{} (\uncertain{général}) : \emph{Dwork}, \emph{Washnitzer}–\emph{Monsky}, plus tard le yoga « cristallin » développé par \emph{Berthelot}, \emph{Illusie}, \emph{Messing}, \add{Mazur} \struck{\ill{}}.
\marginal{au crayon : aussi Hartshorne, Herrera \struck{\ill{}}, Bloch\uncertain{(?)}, Ogus,}
Ici on trouve des théories coh.\ qui ne sont plus à « coeff.\ » de car.\ nulle i.e.\ contenant $\mathbb{Q}$ — et en principe les phénomènes de torsion ne sont pas perdus.
\marginal{au crayon : Du pt de vue Weil, c'était un défaut impardonnable !}

f) Th.\ de \emph{Grothendieck} sur variétés algébriques sur $\mathbb{C}$ (\struck{\ill{}} \add{généralisation} par \emph{Deligne}, \emph{Hartshorne} pour des coeff.\ plus généraux). Ceci donne \uncertain{enfin} (du moins en car.\ 0) \uncertain{la} signific.\ topologique de la coh.\ de De Rham « algébrique » des schémas algébriques.

g) Complexe de De Rham–Sullivan pour espaces topologiques généraux : formes différentielles singulières $C^\infty$ \struck{\ill{}} \struck{\ill{}} ($\mathbb{C}$-algébriques, resp.\ $\mathbb{R}$-algébriques, resp.\ $\mathbb{Q}$-algébrique). Donne la coh.\ à coeff.\ dans $\mathbb{C}$ (resp.\ $\mathbb{R}$, resp.\ $\mathbb{Q}$) (facile). Sullivan montre mieux

\note{En bas à gauche de la page, des inclusions :}
\[
C^\bullet_{\mathrm{DRS}\,\mathbb{R}\text{-alg}} \subset C^\bullet_{\mathrm{DRS}\,C^\infty} \subset C^\bullet_{\mathrm{DRS}},
\qquad C^\bullet_{\mathrm{DRS}\,C^\infty} \supset C^\bullet_{\mathrm{DR}}
\]
\note{la seconde inclusion est écrite verticalement ($\cup$) sous $C^\bullet_{\mathrm{DRS}\,C^\infty}$ ; à sa gauche un symbole biffé.}

\page{3}
\note{Page numérotée (2) par l'auteur.}
que le \struck{$\mathbb{Q}$} type d'homotopie de $X$ est récupéré \add{simplement} si $X$ connexe (sauf erreur) — de façon plus précise, il y a une équivalence \add{donnée par $C^\bullet_{\mathrm{DRS}}$} de cat.\ \struck{\ill{}} \add{faible} entre la cat.\ hom.\ des espaces connexes et simplement connexes (\ill{}), et une certaine catégorie dérivée formée avec les $\mathbb{Q}$-algèbres graduées \struck{\ill{}} associatives et anticommutatives \struck{telle} \add{de} degrés $\geq 0$ telles que $H^0(A) \simeq \mathbb{Q}$ \struck{\ill{}}, $H^1(A) \simeq 0$. Il y a une théorie des modèles minimaux \add{pour de telles algèbres} qui sont une façon \add{\uncertain{algébrique}} très simple de récupérer les $\pi_i \otimes_{\mathbb{Z}} \mathbb{Q}$ en termes d'un tel modèle… (On renvoie à Sullivan)

Mais à nouveau, on perd la torsion !

\textbf{Dém.\ du th.\ d'isomorphisme de Sullivan}
\[ H^q(C^\bullet_{\mathrm{DRS}\,\mathbb{Q}}(X,\mathbb{Q})) \simeq H^q(X,\mathbb{Q}) \]
$\forall n \geq 0$, $[n]$ simplexe type de dim.\ $n$ \quad ($\sum_0^n x_i = 1$, $x_i \geq 0$), $\Delta^{[n]}_{\mathbb{R}}$ sa réalisation géométrique, $E^{[n]}$ l'espace affine qu'il sous-tend (avec une $\mathbb{Q}$-structure) \quad ($\sum_0^n x_i = 1$ dans $\mathbb{R}^{n+1}$),
\[ \mathrm{DRS}^\bullet_{[n]} = C^\bullet_{\mathbb{Q}\text{-DR}}(E^{[n]}) \]
son complexe de De Rham $\mathbb{Q}$-algébrique. C'est contravariant en $[n]$, d'où
\[ \mathrm{DRS}^{\bullet}_{*} = (\mathrm{DRS}^\bullet_{[n]})_{n \geq 0} \]
\uncertain{Algèbre} différentielle graduée \struck{\ill{}} \add{semi-}simpliciale (\uncertain{cot.\ = simplicial}) — à degrés $\geq 0$, anticommutative.

Pour $X$ espace topologique, $S_*(X)$ son complexe ½simplicial. On a
\[ C^\bullet_{\mathrm{DRS}}(X,\mathbb{Q}) \simeq \mathrm{Hom}(S_*(X), \mathrm{DRS}^\bullet_*) \]

\page{4}
\note{Page numérotée (3) par l'auteur.}
i.e.\ a) $C^\bullet_{\mathrm{DRS}}(X,\mathbb{Q})$ dépend de $X$ via $S_*(X)$

b) Sur $(\mathrm{Ss}$ \struck{\ill{}} $(\mathrm{Ens}))_*$, le foncteur $C^\bullet_{\mathrm{DRS}}(-,\mathbb{Q})$ est représentable par $\mathrm{DRS}^\bullet_*$ \quad (dans la cat.\ des \uncertain{gr.} ½simpl.)

Or a) $\mathrm{DRS}^\bullet_*$ est une résolution de $\mathbb{Q}_*$ (lemme de Poincaré algébrique sur les $\mathbb{Q}$-espaces affines $E^{[n]}$)

b) Les composantes $\mathrm{DRS}^i_*$ ($i \geq 0$) de $\mathrm{DRS}^\bullet_*$ sont des objets acycliques du topos $\mathrm{Ss}$ (ce qui revient à dire que leurs $\pi_q(\mathrm{DRS}^i_*)$ sont nuls), ce qui se vérifie facilement. [Ce serait \uncertain{pareil} avec des formes $C^\infty$ à coeff.\ dans $\mathbb{R}$ ou $\mathbb{C}$, ou analytiques réelles, ou analytiques complexes…)

\struck{On aimerait trouver une $\mathbb{Z}$-algèbre diff.\ $C^\bullet_{\mathrm{DR}}(X,\mathbb{Z})$ \ill{} (\ill{} de $X$, ou $S_*(X)$) dans la cat.\ \ill{} \ill{} entière \ill{} $\mathbb{Z}$. Mais il paraît qu'il y a des \ill{} contre ! On ne pourrait y arriver \ill{} \ill{}}
\note{passage de deux lignes et demie biffé par deux longs traits horizontaux ; lecture très incomplète.}

Si on prend $C^\bullet_{\mathbb{Z}\text{-DR}}(E^{[n]})$ (car $E^{[n]}$ a une $\mathbb{Z}$-structure affine) c'est a) qui devient faux : pour intégrer $\int x^n\,dx$, il faut un dénominateur $\dfrac{x^{n+1}}{n+1}$ ! Mais en géom.\ alg.\ on est déjà familiarisé avec une façon de sauter à pieds-joints \add{par dessus} \struck{sur} le canular, consistant à \uncertain{utiliser} les p.d.\ et des polynômes \add{à séries formelles} (\uncertain{puiss.\ div.}). Si $E^{[n]}$ avait une origine sur $\mathbb{Z}$ (i.e.\ pouvait se \uncertain{considérer} comme un $\mathbb{Z}$-module libre de t.f.) — on aurait un complexe de De Rham à p.d.\ Mais \uncertain{par malheur} ça n'en a pas ! Voilà
\[ \mathrm{DRS}^\bullet_{[n]} \simeq \text{\struck{$\mathbb{Q}[X_i, dX_i]$}}\ C^\bullet_{\mathrm{DR}/\mathbb{Q}}\bigl(\mathbb{Q}[(X_i)_{0 \leq i \leq n}]/\textstyle\sum X_i - 1\bigr) \]
\[ = \mathbb{Q}[X_i, dX_i]_{0 \leq i \leq n} \big/ \bigl(\textstyle\sum X_i - 1,\ \sum dX_i\bigr) \simeq \mathbb{Q}[X_0, \ldots, X_{n-1}][dX_0, \ldots, dX_{n-1}] \]

\page{5}
\note{Page numérotée (4) par l'auteur.}
Donc on aurait envie de prendre
\[ \text{\struck{$\mathbb{Q}$}}\ \mathbb{Z}\bigl(\{X_i\}[dX_i]\bigr)_{0 \leq i \leq n} \big/ \bigl(\textstyle\sum X_i - 1,\ \sum dX_i\bigr)_{\mathrm{pd}} \]
où $\{\ \}$ désigne les \uncertain{p.d.\ à p.d.}, mais, \ill{} \ill{}
\[ \simeq \mathbb{Z}\{X_0, \ldots, X_{n-1}\}[dX_0, \ldots, dX_{n-1}] \]
(qui est \struck{la cat} bien une résolution de $\mathbb{Z}$), car si on \uncertain{veut} \uncertain{tenir} \uncertain{de} la relation $\sum_0^n X_i - 1 = 0$,
\[ X_n = 1 - \sum_0^{n-1} X_i \]
(et $dX_n = -\sum_0^{n-1} dX_i$ \uncertain{qui} $\sum_0^n dX_i = 0$)
— on a le « \uncertain{bec} » que $1 - \sum_0^{n-1} X_i$ \uncertain{n'appartient nullement} à l'idéal à p.d.\ donné dans $\mathbb{Z}\{X_0, \ldots, X_{n-1}\}$ (dans ce cas on \uncertain{pourrait} \ill{} \uncertain{envoyer} $\mathbb{Z}\{X_0, \ldots, X_{n-1}\}$ dans $\mathbb{Z}\{X_0, \ldots, X_{n-1}\}$ avec \uncertain{comme} \struck{\ill{}} $\sum_0^n X_i - 1$ dans le \uncertain{noyau}…).

On s'en tire en prenant un anneau de coeff.\ différent de $\mathbb{Z}$, soit $S$, avec un « \uncertain{paramètre} » \struck{\ill{}} $t \in S$ fixé \add{dont \ill{} \ill{}} (i.e.\ $t \in J$, $J$ idéal à p.d.\ de $S$) \struck{\ill{} \ill{} \ill{}} \emph{et remplaçant} \add{les p.d.}\ \struck{\ill{}} les équations $\sum X_i = 1$ de $E^{[n]}$ par l'équation
\[ \sum X_i = t \quad \text{dans } S^{n+1} \]
et définissant
\[ \bigl(C^\bullet_{\mathrm{DRpd}[n]}(S,J,t) = \Bigl(S\{X_i\}_{0 \leq i \leq n}[dX_i]_{0 \leq i \leq n}\Bigr) \Big/ \Bigl(\textstyle\sum X_i - t,\ \sum dX_i\Bigr)_{\mathrm{pd}} \]
(On divise par l'idéal à p.d.\ engendré par $\sum X_i - t$, $\sum_i dX_i$ ; ce qui a un sens

\page{6}
\note{Page numérotée (5) par l'auteur ; la phrase commencée p.~5 se poursuit.}
car dans $S\{X_i\}[dX_i]$ l'idéal formé des formes \add{à p.d.}\ d'augmentation $\in J$ est à p.d.
\[ \Bigl(S\{X_i\}_{0 \leq i \leq n} \big/ (\textstyle\sum X_i - t)_{\mathrm{pd}}\Bigr) \otimes_S \overset{*}{\Lambda}\bigl(S^{[n]}/\mathrm{diag}\,S^{[n]}\bigr) \]
\marginal{dans la marge gauche, tourné : $d(x^{[n]}) = x^{[n-1]}\,dx$ \ill{} \ill{} p.d.}

C'est une $S$-algèbre diff.\ \add{graduée} \uncertain{anticommutative} et de degrés $\geq 0$, \struck{augmentée, \ill{} à p.d.\ sur l'idéal engendré par $\overset{+}{\Lambda}$, $S\{X_i\}^+$ et $J\{X_i\}$} — \uncertain{à p.d.}\ sur l'idéal noyau de l'augmentation
\[ C^\bullet_{\mathrm{DRpd}[n]}(S,J,t) \to S/J \]
et comme telle \uncertain{isomorphe} \ill{} : $S\{X_0, \ldots, X_{n-1}\}[dX_0, \ldots, dX_{n-1}]$, qui est une \emph{résolution de $S$}. \struck{\ill{} a) \ill{}}

\struck{acquis. Mais \uncertain{faute} \ill{} de $S$), \ill{} pour dire \ill{} \ill{} \ill{} $H^q\,C^\bullet$.}
\note{passage encadré et biffé de plusieurs traits obliques.}
\marginal{dans la marge gauche, tourné : \uncertain{i.e.} \ill{}-algèbres gr.\ \uncertain{diff.}\ $S/J$-augmentées, \uncertain{à p.d.\ sur} l'idéal d'augmentation, \uncertain{compatibles} avec la diff.}

Pour $[n]$ variable, on trouve
\[ C^\bullet_{\mathrm{DRpd}*}(S,J,t) = \bigl(C^\bullet_{\mathrm{DRpd}[n]}(S,J,t)\bigr)_{n \geq 0} \]
qui est une rés.\ ½simpliciale (\uncertain{cat.\ = simplicial}) avec aug.\ vers $S/J = k$, et p.d.\ sur l'idéal d'augm., \uncertain{compatible} de $S$. (Elle dépend fonctoriellement de $(S,J,t)$), et elle permet de définir, pour \uncertain{un ensemble ½simplicial} \add{$\mathcal{X}_* \in (\mathrm{Ss})$}
\[ C^\bullet_{\mathrm{DRpd}}(\mathcal{X}_*; S,J,t) = \mathrm{Hom}\bigl(\mathcal{X}_*, C^\bullet_{\mathrm{DRpd}*}(S,J,t)\bigr) \]
foncteur contravariant en $\mathcal{X}_*$, \uncertain{et si $X$ esp.\ top.}
\[ C^\bullet_{\mathrm{DRpd}}(X; \ ) = C^\bullet_{\mathrm{DRpd}}(S_*(X); \ldots) \doteq \mathrm{Hom}\bigl(S_*(X), C^\bullet_{\mathrm{DRpd}*}(\ )\bigr) \]
Mais on \uncertain{ne peut} dire \add{(on a seulement $H^\bullet_{\mathrm{DRpd}}(\mathcal{X}; S,J,t) \to H^\bullet(\mathcal{X}_*,S)$) \uncertain{généralement}} que \ill{} cohomologie \ill{} et \ill{} \ill{} \ill{} \ill{} \uncertain{ne} sera pas isom.\ (\uncertain{pour} $\mathcal{X}_*$ \uncertain{variable}) à $H^\bullet(\mathcal{X}_*,S)$.

\page{7}
\note{Page numérotée (6) par l'auteur.}
\add{Mais,} soit $k$ un anneau commutatif (\uncertain{ou z.\ unit.}) et $T$ une indéterminée, on prendra donc
\[ S = k\{T\}, \quad J = k\{T\}^+ = \mathrm{Ker}\bigl(k\{T\} \to k\bigr), \quad t = T \]
On a
\[ \begin{cases} C^\bullet_{\mathrm{DRpd}[n]}(S,J,t) \simeq \underbrace{k\{T, X_0, \ldots, X_n\}/(\sum X_i - T)_{\mathrm{pd}}}_{\simeq\, k\{X_0, \ldots, X_n\}} \otimes_k \overset{\bullet}{\Lambda}\, k^{[n]}/k \\ S/J \simeq k \end{cases} \]
Soit
\[ \Phi_{k*} = \bigl[[n] \mapsto k^{[n]}\bigr] \supset k_* = \bigl[[n] \to k\bigr] \]
\note{l'inclusion $k_* \subset \Phi_{k*}$ est légendée « sous-module diagonal ».}
\[ \Psi_{k*} = \Phi_{k*}/k_* = \bigl([n] \to k^{[n]}/\underbrace{k}_{\mathrm{diag}}\bigr) \]
On a donc
\[ C^\bullet_{\mathrm{DRpd}*}\bigl(k\{T\}, k\{T\}^+, k\bigr) \simeq \Gamma^\bullet_k \Phi_{k*} \otimes_k \overset{\bullet}{\Lambda}\, \Psi_{k*} \]
\note{le membre de gauche est noté, par un signe d'égalité vertical légendé « déf », $C^{\bullet\bullet}_{\mathrm{DRpd}*}(k)$ ; formule encadrée :}
\[ C^{\bullet\bullet}_{\mathrm{DRpd}*}(k) \simeq \Gamma^\bullet_k \Phi_{k*} \otimes_k \overset{\bullet}{\Lambda}\, \Psi_{k*} \]

\emph{Structure} : \struck{$k\{T\}$-} \struck{\ill{}} Algèbre \emph{bigraduée} (deg.\ complém.\ \emph{et} deg.\ cohom., \emph{d'où} degré total) \add{unit.}\ assoc.\ \emph{alternée} (anticommutative + carré des éléments de degré impair nuls), augmentée vers $k$, à p.d.\ sur l'idéal d'augm., avec diff.\ de bidegré $-1,+1$ (donc degré total 0) compatible avec p.d.
\marginal{dans la marge gauche, en oblique : la bigraduation \uncertain{permet} de \ill{} \ill{} gr.\ \struck{\uncertain{anticomm.}} \ill{} \ill{} (\uncertain{ici de degré 2})}

\page{8}
\note{Page numérotée (7) par l'auteur.}
Ces structures sont héritées par les
\[ C^{\bullet\bullet}_{\mathrm{DRpd}}(\mathcal{X}_*, k) \overset{\mathrm{déf}}{=} C^\bullet_{\mathrm{DRpd}}\bigl(\mathcal{X}_*; k\{T\}, k\{T\}^+, T\bigr) = \mathrm{Hom}\bigl(\mathcal{X}_*, C^{\bullet\bullet}_{\mathrm{DRpd}*}(k)\bigr) \]
\note{formule encadrée.}
et dépendent de façon contravariante de $\mathcal{X}_*$ (covariante de $k$).

$\Phi_{k*}$ est \add{un $k$-Module ½simplicial} \struck{\ill{}} homotope à 0, donc $\Gamma^p_k(\Phi_{k*}) \otimes \overset{q}{\Lambda}\,\Psi_{k*}$ \struck{\ill{}} est homotope à 0 si $p \neq 0$.
\note{à gauche, une ligne inachevée : $\Gamma^p_k(0) \otimes \overset{q}{\Lambda}\,\Psi_{k*}$, donc}

Donc $C^{p,q}_{\mathrm{DRpd}*}(k)$ est homot.\ à 0 (donc \add{acycliques} \struck{flasques}) si $p \neq 0$. En degré total donné $n$, on a
\[ 0 \to C^{n,0}_* \to C^{n-1,1}_* \to C^{n-2,2}_* \to \cdots \to C^{1,n-1}_* \to C^{0,n}_* \to 0, \qquad k_* \to C^{n,0}_* \]
\note{l'augmentation $k_* \to C^{n,0}_*$ est une flèche verticale montante ; à l'extrémité droite de la suite, la flèche entre $C^{1,n-1}_*$ et $C^{0,n}_*$ semble tracée vers la gauche, et quelques signes tombent dans la marge : lecture \uncertain{incertaine} de cette extrémité.}

i.e.\ on trouve une résolution \add{de longueur $n$} \struck{flasque} de $k_*$ par des \struck{\ill{}} $k$-Modules ½simpliciaux qui sont acycliques sauf le dernier — donc ce qui précède le dernier comme un tronqué \add{degré} \struck{l'indice} $n$ d'une résolution flasque de $k_*$ — d'où le calcul de la cohomologie de ses sections sur un $\mathcal{X}_*$ : c'est la coh.\ de $\mathcal{X}_*$ tronquée en degré $n$ :
\[ H^{p,q}_{\mathrm{DRpd}}(\mathcal{X}_*, k) \simeq \begin{cases} H^q(\mathcal{X}_*, k) & \text{si } q \leq p+q \text{ i.e.\ } p \geq 0 \\ 0 & \text{si } p < 0 \end{cases} \]
\note{la condition « $q \leq p+q$ » est une lecture \uncertain{incertaine}.}

\page{9}
\note{Page numérotée (8) par l'auteur.}
Mais quelle est \add{(pour $q$ fixé)} la structure de $k\{T\}$-module de $H^{\bullet q}_{\mathrm{DRpd}}(\mathcal{X}_*, k)$ ? On \struck{voit que} \struck{\ill{}} pour le degré total (égal au degré cohomologique \emph{plus} $q$) on trouve $H^0(\mathcal{X}_*, k) \otimes_k k\{T\}$ tronqué en degré $\geq q$, donc
\[ H^{\bullet q}_{\mathrm{DRpd}}(\mathcal{X}_*, k) \simeq \tau_q\bigl(H^0(\mathcal{X}_*, k) \otimes_k k\{T\}\bigr)[q] \]
\note{l'auteur renvoie par une flèche sur le $[q]$ : « (translation des degrés, $p \mapsto -q$) ».}

Si on réindexe les bidegrés par le complexe (degré total, degré cohomique) \uncertain{posant}
\[ H'^{\,n,q}_{\mathrm{DRpd}}(\mathcal{X}_*, k) = H^{\overset{p}{\overbrace{n-q}},\,q}_{\mathrm{DRpd}}(\mathcal{X}_*, k) \]
(dans les conditions de degré $p, q \geq 0$ deviennent $n \geq q \geq 0$), l'op.\ diff.\ est de bidegré $(0,1)$ donc c'est un \uncertain{homomorphisme} (\emph{homogène}) de $S$-modules gradués), on trouve
\[ H'^{\,\bullet,q}_{\mathrm{DRpd}}(\mathcal{X}_*, k) \simeq \tau_q\bigl(H^q(\mathcal{X}_*, k) \otimes_k k\{T\}\bigr) \]
Ces isom.\ sont compatibles avec les structures multiplicatives (\uncertain{c'est bien} \uncertain{évidemment} \ill{}, \uncertain{voir} $\mathcal{X}_*$, $k$…). Donc \add{a priori} \uncertain{on} \uncertain{récupère} \add{(\uncertain{via} $H'^{\bullet\bullet}_{\mathrm{DRpd}}(\mathcal{X}_*, k)$)} \uncertain{par} les $k$-modules $H^q(\mathcal{X}_*, k)$ et leurs \add{cup-}accouplements, mais les $k$-modules \add{gradués} $\tau_q\bigl(H^q(\mathcal{X}_*, k) \otimes_k k\{T\}\bigr)$ et leurs accouplements déduits des cup dans $H^\bullet(\mathcal{X}_*, k)$.

\page{10}
\note{Page numérotée (9) par l'auteur.}
Ce qui donne une \uncertain{espérance} que le complexe de DR à p.d.\ de $\mathcal{X}_*$ \add{à coeff.\ dans $k = \mathbb{Z}$ disons} (\uncertain{comme} dans le cas $k = \mathbb{Q}$) qu'il contient) donne \uncertain{bien} \uncertain{une} information homotopique précise sur $\mathcal{X}_*$, c'est l'observation que si on fait, pour un \struck{module} \add{gradué} $k$-module $M$ quelconque, on récupère $M$ à isom.\ \uncertain{près} connaissant \uncertain{pris} par le connaissance d'un quelconque des tronqués $\tau_q(M \otimes_k S)$ (quelque \add{\uncertain{fixe}} \uncertain{soit} $q$ \add{\uncertain{tels} \uncertain{que} \ill{} au-dessus de $M$}) \uncertain{passant par ses $\rho$}…) (et de \uncertain{même} \uncertain{l'}accouplement $M \otimes N \to P$ $\to$ connaissant, pour $q$ assez grand, l'accouplement correspondant $\tau_q(M_S) \otimes \tau_q(N_S) \to \tau_q(P_S)$.

Donc une $k$-algèbre gr.\ $H^\bullet$ à degrés $\geq 0$ \struck{\ill{}} est connue à isom.\ canonique près quand on connaît la $S$-algèbre bigraduée \add{de degré $q$} des « extinctions » \add{$\tau_q$} \uncertain{sur} les $\tau_q(H^q \otimes_k S)$…. \struck{Ainsi} Cette « théorie des ombres » étant \add{\uncertain{admise}} acquise, \uncertain{on} \uncertain{voit} que la connaissance de $C^{\bullet\bullet}_{\mathrm{DRpd}}(\mathcal{X}_*, k)$ (\uncertain{implique} \add{\uncertain{comme} \uncertain{les} \uncertain{plus} $S$-algèbre diff.\ bigraduée}) implique celle de l'algèbre graduée $H^\bullet(\mathcal{X}_*, k)$ — et \uncertain{en} \uncertain{affinant} un \uncertain{peu}, \uncertain{on} \add{\uncertain{voit} \uncertain{que}} \struck{\ill{}} \uncertain{pour} $\mathrm{R}\Gamma^\bullet(\mathcal{X}_*, k)$ comme objet de la cat.\ dérivée $D(k)$, \struck{\ill{}} l'accouplement

\page{11}
\note{Page numérotée (10) par l'auteur ; la phrase de la p.~10 se poursuit.}
\uncertain{venant} de cup-produits
\[ \mathbb{R}\Gamma(\mathcal{X}_*, k) \overset{L}{\otimes} \mathbb{R}\Gamma(\mathcal{X}_*, k) \to \mathbb{R}\Gamma(\mathcal{X}_*, k) \]

\textbf{Remarques et Pbs}

a) Si $k$ est une $\mathbb{Q}$-algèbre, il se \uncertain{trouve} que $C^\bullet_{\mathrm{DRS}}(\mathcal{X}_*, k)$ est \uncertain{défini}, et se reconstruit à partir de $C^{\bullet\bullet}_{\mathrm{DRpd}}(\mathcal{X}_*, k)$ \struck{\ill{}} par la formule
\[ C^\bullet_{\mathrm{DRS}}(\mathcal{X}_*, k) \simeq C^{\bullet\bullet}_{\mathrm{DRpd}}(\mathcal{X}_*, k)/(T-1) \]
Plus généralement, quel que soit $k$, on a
\[ C^\bullet_{\mathrm{DRS}}(\mathcal{X}_*, \underbrace{k \otimes_{\mathbb{Z}} \mathbb{Q}}_{k_{\mathbb{Q}}}) \simeq C^{\bullet\bullet}_{\mathrm{DRpd}}(\mathcal{X}_*, k)/(T-1) \]
[et plus général\uncertain{ement}, si $(S,J,t)$ comme ci-dessus
\[ C^\bullet_{\mathrm{DRpd}}(\mathcal{X}_*; S,J,t)/(t-1) \simeq C^\bullet_{\mathrm{DRS}}\bigl(\mathcal{X}_*, S_{\mathbb{Q}}/(t-1)\bigr) \]
\note{le membre de droite est en partie surchargé ; un « ? » en marge gauche.}
]

b) \uncertain{J'ai} négligé \uncertain{jusqu'ici} la structure à p.d.\ sur l'idéal d'augmentation de
\[ C^{\bullet\bullet}_{\mathrm{DRpd}}(\mathcal{X}_*, k) \to H^0(\mathcal{X}_*, k) = \mathrm{Hom}(\mathcal{X}_*, k_*) \]
\uncertain{qui} est importante. \struck{\uncertain{Je crois} \uncertain{que} \ill{} \uncertain{par}} \uncertain{L'idée} de la question, pour $\mathcal{X}_*$ simplement \uncertain{connexe}, si la connaissance de $C^{\bullet\bullet}_{\mathrm{DRpd}}(\mathcal{X}_*, k)$ \emph{avec toutes ses structures} (y compris celles des p.d.) implique \uncertain{par} la connaissance du type d'homotopie de $\mathcal{X}_*$. On peut
\marginal{dans la marge gauche, en oblique : avec conditions de finitude sur $\mathcal{X}_*$ \uncertain{et} \uncertain{sur} \uncertain{les} $\pi_i(\mathcal{X}_*)$ \ill{}}
\note{la phrase se poursuit au-delà de la p.~11.}

\page{12}
\note{Page numérotée (11) par l'auteur ; la phrase de la p.~11 se poursuit.}
plus précisément \uncertain{appelons} \emph{complexe de De Rham à p.d.\ virtuel} \add{sur $k$} \struck{\uncertain{un} complexe de} $k$-bialgèbre diff.\ \struck{\ill{}} augmentée, ass.\ unit.\ alternée, à diff.\ de bidegré $-1,+1$ \struck{(\uncertain{ou} $0,+1$ \uncertain{au} \uncertain{choix})} à bidegrés $\geq 0$, avec p.d.\ sur l'idéal d'augm., compatibles avec la diff.\ ($d(x^{[n]}) = x^{[n-1]}\,dx$), et telle que les $H^{\bullet q}(C^{\bullet\bullet})[-q]$ \struck{\ill{}} \uncertain{soient} des $q$-ombres (auquel cas on \uncertain{peut} récupérer à partir de $C^{\bullet\bullet}$ un élément de $D(k)$ avec structure multiplicative \uncertain{ass.\ unit.}\ commutative….), \uncertain{posons} : une « catégorie dérivée » des \uncertain{ces} complexes en inversant les flèches qui sont des \uncertain{q-isom.}, d'où par $C^{\bullet\bullet}_{\mathrm{DRpd}}(\mathcal{X}_*, k)$ un foncteur
\[ (\mathrm{Hot}) \to (\mathrm{DRpd}) \]
\note{sous $(\mathrm{Hot})$ : « type d'hom. » ; sous $(\mathrm{DRpd})$ : « cat.\ dérivée des complexes DR pd ».}
et on peut se demander si sa restriction aux espaces connexes et simplement connexes avec des $H^i(X,\mathbb{Z})$ de t.f.\ (\uncertain{pour} $k = \mathbb{Z}$) \struck{\ill{}} \add{induit} une équivalence avec les complexes de DRpd sur $\mathbb{Z}$ tels que $H^0(C^{\bullet\bullet}) \simeq k$, $H^1(C^{\bullet\bullet}) = 0$, les $H^i(C^{\bullet\bullet})$ de t.f.\ sur $\mathbb{Z}$… !

\page{13}
\note{Page numérotée (12) par l'auteur.}
c) J'ai un peu réfléchi \struck{\ill{}} si on peut reconstruire les op.\ coh.\ \struck{\ill{}} (type Steenrod ou Whitney) dans la coh.\ de $\mathcal{X}_*$ par la connaissance de $C^{\bullet\bullet}_{\mathrm{DRpd}}(\mathcal{X}_*, -)$, et n'ai que des résultats \add{partiels} négatifs qui montrent qu'avec des \uncertain{autres} \struck{\ill{}} multiplicatifs de $C^{\bullet\bullet}_{\mathrm{DRpd}*}(k)$ on ne trouve rien d'intéressant.

d) Il faudrait \uncertain{donc} \uncertain{avant} \uncertain{étudier} des modèles minimaux à la Sullivan, pour essayer entre autres d'exprimer \uncertain{les} \uncertain{invariants} \uncertain{du} type d'homotopie de $\mathcal{X}_*$ en termes de $C^{\bullet\bullet}_{\mathrm{DRpd}*}(\mathcal{X}_*, \mathbb{Z})$ (\uncertain{au} \uncertain{moins} \uncertain{si} $\mathcal{X}_*$ \struck{simpl.} connexe avec cond.\ de finitude…) Je n'ai rien fait dans cette direction, je n'ai \uncertain{même} pas développé une formule de Künneth pour $C^{\bullet\bullet}_{\mathrm{DRpd}}$ d'un produit $\mathcal{X}_* \times \mathcal{Y}_*$ — il y a des difficultés techniques dues au fait qu'on ne peut sans doute supposer $C^{\bullet\bullet}_{\mathrm{DRpd}}(\mathcal{X}_* \text{ ou } \mathcal{Y}_*, k)$ plats \struck{\ill{}} sur $k\{T\}$.

\page{14}
\note{Page numérotée (13) par l'auteur.}
e) Le complexe de chaînes $C_\bullet(\mathcal{X}_*,$ \struck{$k$} $\mathbb{Z})$ permet de reconstruire tous les complexes de cochaînes $C^\bullet(\mathcal{X}_*, k')$, pour $k'$ \struck{\ill{}} \add{$k$-alg.}\ \uncertain{quelconque}, par
\[ C^\bullet(\mathcal{X}_*, k) \simeq \mathrm{Hom}^\bullet_{\mathbb{Z}}\bigl(C_\bullet(\mathcal{X}_*, \mathbb{Z}), k\bigr) \]
(L'objet « le plus fin » est $C_\bullet(\mathcal{X}_*, \mathbb{Z})$, on a $C_\bullet(\mathcal{X}_*, k) \simeq C_\bullet(\mathcal{X}_*, \mathbb{Z}) \otimes_{\mathbb{Z}} k$.) Il est possible de \struck{\ill{}} de définir une cobigèbre \add{diff.}\ $C^{\mathrm{DRpd}}_{\bullet\bullet}(\mathcal{X}_*,$ \struck{$k$} $k)$ (à \uncertain{comultipl.}\ \uncertain{div.}, \add{$k$-coaugmentation}), telle que l'on ait, pour toute $k$-alg.\ $k'$
\[ C^{\bullet\bullet}_{\mathrm{DRpd}}(\mathcal{X}_*, k') \simeq \mathrm{Hom}_k\bigl(C^{\mathrm{DRpd}}_{\bullet\bullet}(\mathcal{X}_*, k), k'\bigr) \]
(compatible avec toutes les structures). On aura d'ailleurs
\[ C^{\mathrm{DRpd}}_{\bullet\bullet}(\mathcal{X}_*, k') \simeq C^{\mathrm{DRpd}}_{\bullet\bullet}(\mathcal{X}_*, k) \otimes_k k' \]
L'objet le plus fin est $C^{\mathrm{DRpd}}_{\bullet\bullet}(\mathcal{X}_*, \mathbb{Z})$. C'est lui qu'il conviendrait de considérer (au lieu de son « dual » $C^{\bullet\bullet}_{\mathrm{DRpd}}(\mathcal{X}_*, \mathbb{Z})$) si on veut aborder b) c) d) sans conditions de finitude.

Notons que [tout \uncertain{comme} $C^{\bullet\bullet}_{\mathrm{DRpd}}(\mathcal{X}_*, k)$, pour $\mathcal{X}_*$ variable transforme $\varinjlim$ quelconques en $\varprojlim$] $C^{\mathrm{DRpd}}_{\bullet\bullet}(\mathcal{X}_*, \mathbb{Z})$ \struck{transforme} \add{\uncertain{commute}} aux $\varinjlim$ quelconques (\uncertain{pas} \uncertain{seulement} \uncertain{filtrantes}).

\page{15}
\note{Page numérotée 14 par l'auteur.}
\marginal{en oblique, en haut à gauche, reliée par un trait à la lettre de l'alinéa : \uncertain{voir} (f) ci-dessous}
\struck{f)} g) \emph{Faisceautisation}. Il y a une définition évidente de complexes de De Rham à p.d.\ sur $k$ si $k$ est un Anneau comm.\ d'un topos. On aimerait, p.\ ex.\ en comparant un tel complexe à un autre \uncertain{q.-isom.}\ dont les composantes soient flasques (mais y en a-t-il \uncertain{toujours} ?) définir des \uncertain{opér.}\ $\mathrm{R}f_*$ dans des cat.\ dérivées convenables pour de tels complexes, quand $f \colon \mathcal{X} \to \mathcal{Y}$ est un morphisme de topos \add{(\uncertain{supposés} \uncertain{peut-être} de \uncertain{dim.}\ coh.\ finie…)} si $k$ est une $\mathbb{Q}$-Algèbre, le \uncertain{même} Pb se \uncertain{rencontre} d'ailleurs déjà pour les complexes du type De Rham–Sullivan, c'est le Pb \uncertain{ouvert} — et posé par Deligne — quand $\mathcal{X} = (\mathrm{Ens})^\Delta =$ topos des ens.\ ½simpliciaux, $\mathcal{Y} =$ topos ponctuel. Mais sauf erreur (si souvenirs \uncertain{sont} exacts) il y a un topos qui marche, pour les espaces \add{topologiques} \struck{\uncertain{paracompacts}}….

\page{16}
\note{Page numérotée (15) par l'auteur.}
\struck{g)} h) On peut associer à un espace topologique $X$ \add{\uncertain{ou} à} un ens.\ ½simplicial $\mathcal{X}_*$ \struck{\ill{}} \add{\uncertain{des} invariants} \uncertain{alg.}\ « linéaires » plus fins que ceux \uncertain{fournis} par le $C^{\bullet\bullet}_{\mathrm{DRpd}}$ (ou \uncertain{mieux} par $C^{\mathrm{DRpd}}_{\bullet\bullet}$, \uncertain{en} \uncertain{dualisant}…). [P.\ ex.\ on peut observer que
\[ C^{\bullet\bullet}_{\mathrm{DRpd}*}(k) \simeq \Gamma^\bullet \Phi_{*k} \otimes_k \overset{\bullet}{\Lambda}\, \Psi_{*k} \]
se déduit de $\mathcal{D}^{\bullet\bullet}_{*(k)} \overset{\mathrm{déf}}{=} \Gamma^\bullet \Phi_{*k} \otimes \overset{\bullet}{\Lambda}\, \Phi_{*k}$ (qui est une \add{$k$-\uncertain{augm.}} alg.\ diff.\ à p.d.\ qui est une \emph{résolution de $k$}) et de $T \in \Gamma(\mathcal{D}^{10}_{*k})$ comme \uncertain{noyau} de la multiplication par $dT$ (i.e.\ on divise par l'idéal \struck{engendré} \add{\uncertain{annulateur}} par $dT$), en \uncertain{chaque} bidegré \struck{\ill{}} $(p,q)$ \add{donné}
\[ \mathcal{C}^{p,q}_{*k} \simeq \mathrm{Ker}\bigl(\mathcal{D}^{p,q+1}_* \to \mathcal{D}^{p,q+2}_*\bigr) \]
\note{sous la flèche : « produit par $dT$ » ; les exposants $q+1$, $q+2$ sont en partie surchargés.}
d'où
\[ C^{p,q}_{\mathrm{DRpd}}(\mathcal{X}_*, k) \simeq \mathrm{Ker}\bigl(\mathcal{D}^{p,q+1}(\mathcal{X}_*, k) \to \mathcal{D}^{p,q+2}(\mathcal{X}_*, k)\bigr) \]
et on peut considérer les structures \add{\uncertain{diverses}} sur $C^{\bullet\bullet}_{\mathrm{DRpd}}(\mathcal{X}_*, k)$ comme déduites de certaines structures (\uncertain{compliquées}…) sur $\mathcal{D}^{\bullet\bullet}(\mathcal{X}_*, k)$.
On peut aussi définir « \emph{La} structure multiplicative à p.d.\ \add{\uncertain{cohomologique}} du type d'homotopie $\mathcal{X}_*$ », en un sens \uncertain{convenable}, qui permettra de reconstituer \struck{\ill{}} \add{\uncertain{aussi} \uncertain{bien}} $C^{\bullet\bullet}_{\mathrm{DRpd}}(\mathcal{X}_*, k)$ que $\mathcal{D}^{\bullet\bullet}_k(\mathcal{X}_*, k)$ (ou les complexes de De Rham–Sullivan). J'ignore si la
\note{la phrase se poursuit au-delà de la p.~16 ; le crochet ouvert « [P.\ ex. » n'est pas refermé sur la page.}

\page{17}
\note{Page numérotée (16) par l'auteur ; la phrase de la p.~16 se poursuit et s'achève ici, sur la dernière page de la série numérotée.}
complexe de De Rham à p.d.\ \add{(ou même $\mathcal{D}^{\bullet\bullet}(\mathcal{X}_*, \mathbb{Z})$)} suffit déjà pour récupérer toute la structure multipl.\ à p.d.\ coh.\ de $\mathcal{X}_*$ \add{(sous réserve de conditions de finitude bien sûr)}. Dans le cas contraire, ce serait cette dernière \add{alg.\ « linéaire »} qui serait le candidat \struck{\uncertain{naturel}} pour exprimer le type d'homotopie $\mathcal{X}_*$ (du moins si $\mathcal{X}_*$ \struck{simplement} \add{\uncertain{connexe}} connexe, et en passant à une cat.\ dérivée convenable bien sûr).

\section{Lemmes sur les coeff.\ binomiaux}
\note{titre encadré, en haut à droite de la p.~19, de la main de l'auteur. Les pp.~19--20 ne portent pas de numéro de l'auteur ; la p.~18 est une feuille de couleur vierge.}

\page{19}
\textbf{Lemme 1.} \emph{Soient $p$ un premier, $i, j \in \mathbb{N}$, avec des développements $p$-adiques
\[ i = \sum_{0 \leq \varepsilon \leq e} u_\varepsilon p^\varepsilon, \qquad j = \sum_{0 \leq \varepsilon \leq e} v_\varepsilon p^\varepsilon, \qquad 0 \leq u_\varepsilon, v_\varepsilon \leq p-1, \]
soit
\[ c_{ij} = \frac{(i+j)!}{i!\,j!} \in \mathbb{Z} \qquad \text{(coeff.\ binomial)} \]
alors on a
\[ c_{ij} \not\equiv 0 \ (p) \quad \text{ssi} \quad u_\varepsilon + v_\varepsilon \leq p-1 \quad \forall \varepsilon \ (0 \leq \varepsilon \leq \nu) \]
et \struck{\ill{}} on a en tous cas
\[ c_{ij} \equiv \prod_{0 \leq \varepsilon \leq \nu} c_{u_\varepsilon, v_\varepsilon} \quad (p) \]}
\note{les bornes sont écrites $e$ dans les sommes, $\nu$ dans la condition et le produit.}

\marginal{en marge gauche de l'alinéa suivant : Cor 1}
On écrit d'abord $i = u + pr$, $j = v + ps$, $0 \leq u, v \leq p-1$, $r, s \geq 0$, et on montre que
\[ (*) \quad \begin{cases} c_{ij} \not\equiv 0 \ (p) \quad \text{ssi} \quad u + v \leq p-1 \text{ et } c_{r,s} \not\equiv 0 \ (p) \\ c_{ij} \equiv c_{u,v}\, c_{r,s} \ (p) \end{cases} \]
en étudiant le développement, dans $\mathbb{F}_p[V, Y]$, de $(X+Y)^{i+j}$. \struck{\ill{}} On a
\[ i + j = p\underbrace{(r+s)}_{t} + \underbrace{u+v}_{w} \]
\note{l'anneau est écrit $\mathbb{F}_p[V, Y]$ alors que le calcul porte sur $X$, $Y$ ; lecture du « $V$ » \uncertain{incertaine}.}
Si $w \leq p-1$ on écrit
\[ (X+Y)^{i+j} = \bigl((X+Y)^p\bigr)^t (X+Y)^w = (X^p + Y^p)^t (X+Y)^w \]
\[ = \sum_{\substack{0 \leq \alpha \leq w \\ 0 \leq \beta \leq t}} \binom{w}{\alpha} \binom{t}{\beta} X^{\alpha + p\beta}\, Y^{(w-\alpha) + p(t-\beta)}, \]
où tous les monômes écrits dans la somme ont \uncertain{des} exposants distincts. On a donc comme coeff.\ de $X^i Y^j$ celui correspondant au terme $i = \alpha + p\beta$ i.e.\ $\alpha = u$, $\beta = r$ i.e.
\[ c_{i,j} \equiv \binom{w}{u} \binom{t}{r} = c_{u,v}\, c_{r,s} \quad (p) \]
et il est immédiat que $c_{u,v} \not\equiv 0$ $(p)$, car $c_{u,v} = \dfrac{(u+v)!}{u!\,v!}$, $0 \leq u, v, u+v \leq p-1$. Donc $c_{ij} \not\equiv 0$ $(p)$ $\Leftrightarrow$ $c_{v,s} \not\equiv 0$ $(p)$.
\note{l'indice est lu $c_{v,s}$ ; on attendrait $c_{r,s}$.}

Il reste à voir que si $w \geq p$, \add{alors} $c_{ij} \equiv 0$ $(p)$. \struck{Soit $w = p + w$} Soit $w = p + w'$, $0 \leq w' \leq p-2$, donc $i + j = (r+s+1)p + w' = \underbrace{(t+1)}_{t'}p + w'$,
\[ (X+Y)^{i+j} = (X^p + Y^p)^{t'} (X+Y)^{w'} = \sum_{\substack{0 \leq \alpha \leq w' \\ 0 \leq \beta \leq t'}} \binom{w'}{\alpha} \binom{t'}{\beta} X^{\alpha + p\beta}\, Y^{(w'-\alpha) + p(t'-\beta)} . \]
Terme en $X^i Y^j$ correspond\supplied{ant} à $\alpha = u$, $\beta = r$, il faudrait $u = \alpha \leq w' = u + v - p$ donc $v \geq p$, ce que \uncertain{n'est} \uncertain{pas}, donc il \uncertain{n'y} a pas de

\page{20}
tels termes. Donc $c_{ij} \equiv 0$ $(p)$ dans ce cas, ce qui prouve (remplaçant $i, j$ par $u, v$) que $c_{u,v} \equiv 0$ $(p)$ donc encore $c_{ij} \equiv c_{u,v}\, c_{r,s}$ $(p)$.

Le lemme 1 résulte du corollaire par une récurrence immédiate sur la longueur $e$ des développements $p$-adiques écrits de $u$ et de $v$.

\textbf{Cor.\ 2.} \emph{Soit $N$ un entier $\geq 0$ tel que $N + 1 \equiv 0$ $(p)$, posons $N + 1 = (N'+1)p$ donc $N = (N'+1)p - 1 = N'p + (p-1)$ (de sorte qu'à décalage de 1 près, les \uncertain{dév.}\ de $N+1$ et $N'+1$ sont $=$, \uncertain{diff.}\ du développement $p$-adique, ainsi que $N$ et $N'$ (quand on fait abstraction du premier coeff.\ $p-1$ de $N$)). Soit alors $\varepsilon \geq 1$, on a}
\[ c_{N+1,\, p^\varepsilon - 1} \equiv c_{N'+1,\, p^{\varepsilon-1} - 1} \quad \begin{cases} \equiv 1 \ (p) & \text{si } N'+1 \equiv 0 \ (p^{\varepsilon-1}) \\ \equiv 0 \ (p) & \text{si } N'+1 \not\equiv 0 \ (p^{\varepsilon-1}) \end{cases} \]
\[ c_{N,\, p^\varepsilon} \equiv c_{N',\, p^{\varepsilon-1}} \quad \Bigl\{ \equiv v_{\varepsilon-1} + 1 \ (p) \quad \Bigl(\text{si } N' = \sum v_\alpha p^\alpha\Bigr) \]
\note{dans les deux lignes, des signes intermédiaires sont biffés de hachures serrées (on devine des « $(p)$ » et « $\not\equiv 0$ ») ; non transcrits.}

\emph{Dém.} Le développement $p$-adique de $p^\varepsilon - 1$ est
\[ p^\varepsilon - 1 = (p-1) + (p-1)p + \cdots + (p-1)p^{\varepsilon-1} \]
de même
\[ p^{\varepsilon-1} - 1 = (p-1) + \cdots + (p-1)p^{\varepsilon-2} \quad \text{si } \varepsilon \geq 2 \qquad (= 0 \text{ si } \varepsilon = 1) \]
Posant
\[ \left.\begin{aligned} N'+1 &= u_0 + u_1 p + \cdots + u_\nu p^\nu \\ N' &= v_0 + v_1 p + \cdots + v_\nu p^\nu \end{aligned}\right\} \quad (\nu \geq \varepsilon - 1) \]
d'où
\[ \begin{aligned} N+1 &= 0 + u_0 p + u_1 p^2 + \cdots + u_\nu p^{\nu+1} \\ N &= (p-1) + v_0 p + v_1 p^2 + \cdots + v_\nu p^{\nu+1} \end{aligned} \]
on trouve
\[ \begin{cases} c_{N+1,\, p^\varepsilon - 1} \equiv c_{u_0, p-1}\, c_{u_1, p-1} \cdots c_{u_{\varepsilon-2}, p-1} \ (p) & (\equiv 1 \ (p) \text{ si } \varepsilon = 1) \\ c_{N'+1,\, p^{\varepsilon-1} - 1} \equiv c_{u_0, p-1}\, c_{u_1, p-1} \cdots c_{u_{\varepsilon-2}, p-1} \ (p) & (\equiv 1 \ (p) \text{ si } \varepsilon = 1) \end{cases} \]
\[ \begin{cases} c_{N,\, p^\varepsilon} \equiv c_{v_{\varepsilon-1}, 1} \equiv v_{\varepsilon-1} + 1 \\ c_{N',\, p^{\varepsilon-1}} \equiv c_{v_{\varepsilon-1}, 1} \equiv v_{\varepsilon-1} + 1 \end{cases} \]
\note{dans la première accolade, le produit est écrit jusqu'à $u_{\varepsilon-2}$ dans les deux lignes.}

\textbf{Lemme 2} \struck{\ill{}}. \emph{Soit $\mathbb{Z}_p$ le localisé de $\mathbb{Z}$ en $p$, $E$ un $\mathbb{Z}_p$-module, $N \in \mathbb{N}$, $e \geq 0$ tel que $p^e \geq N + 1$ \struck{\ill{}}, $I$ l'ens.\ des entiers de la forme $N$ et $N + p^\varepsilon$, $0 \leq \varepsilon \leq e$, $(x_i)_{i \in I}$ une famille d'éléments de $E$, tel qu'on ait
\[ \binom{j}{i}(x_j - x_i) = 0 \qquad i = N, N+1,\ i < j \in I \quad \Bigl(\text{donc } j = N + p^\varepsilon,\ 0 \leq \varepsilon \leq e\Bigr) \]}
\marginal{dans la marge gauche, tourné : Rmq si $E$ est un $\mathbb{F}_p$-module, on ne peut conclure que $x_i = x_N$ $\forall i \in I$ (ex.\ $p = 2$, $N = 2$, $e = 2$, donc \uncertain{on} \uncertain{a} $x_2, x_3, x_4, x_6$ on a $x_2 = x_3 = x_6$ mais \ill{})}
\note{l'énoncé du lemme 2 se poursuit au-delà de la p.~20 ; le batch s'arrête ici.}

\end{document}
